\documentclass[11pt,twoside]{article}

\usepackage[
  a4paper,
  margin=1.5cm,
  includehead, includefoot
]{geometry}

\usepackage[utf8]{inputenc}
\usepackage[brazil,portuges]{babel}
\usepackage{csquotes}

% \usepackage[
%   backend=biber,
%   style=abnt,
%   language=auto,
% ]{bib>}

%   \addbibresource{../refs/referencias.bib}

\usepackage{soul}
\usepackage{xcolor}
\usepackage{framed}

\usepackage{tikz}

  \usetikzlibrary{shapes.geometric, positioning, decorations.pathreplacing, backgrounds, calc, arrows.meta, bending}

\usepackage{kpfonts}
\usepackage[T1]{fontenc}

\usepackage{tabularx}

  \newcolumntype{C}{>{\centering\arraybackslash}X}

\usepackage{graphicx}
\usepackage{multirow}
\usepackage{multicol}
\usepackage{hhline}
\usepackage{booktabs}
\usepackage[inline]{enumitem}
\usepackage{xfrac}
\usepackage{mathtools}

  \newcommand{\Binary}{\ensuremath{\mathbb{B}}}
  \newcommand{\Natural}{\ensuremath{\mathbb{N}}}
  \newcommand{\Integer}{\ensuremath{\mathbb{Z}}}
  \newcommand{\Rational}{\ensuremath{\mathbb{Q}}}
  \newcommand{\Real}{\ensuremath{\mathbb{R}}}

  \DeclarePairedDelimiter{\floor}{\lfloor}{\rfloor}
  \DeclarePairedDelimiter{\ceil}{\lceil}{\rceil}
  \DeclarePairedDelimiter{\abs}{\lvert}{\rvert}

\usepackage{amsthm}

  \newtheorem{proposition}{Proposição}[section]
  \theoremstyle{definition}
  \newtheorem{definition}{Definição}[section]

\usepackage{multicol}

  \newcommand{\finishcols}{\vfill\vphantom{.}}

\usepackage[portuguese,algoruled,vlined,noend]{algorithm2e}

  \SetKwInput{Algoritmo}{Algoritmo}
  \SetKw{Retorne}{retorne}
  \SetKwFor{Para}{para}{faça}{fimpara}
  \SetKw{Escreva}{escreva}
  \DontPrintSemicolon

\usepackage{hyperref}
\usepackage[brazil]{varioref}
\usepackage[brazilian]{cleveref}

  \def\disciplinacodigotext{QXD0115}
  \def\disciplinanometext{Estrutura de Dados Avançada}
  \def\semestretext{2022.1}
  \def\provatype{Notas de Aula}
  \def\provatitle{Rotações e Balanceamento de AVLs}
  \def\provatext{\provatype{} \ifx\provatype\empty {} \else {--} \fi \provatitle{}}
  \def\provacodtext{}
  \def\professorabbrv{Prof.}
  \def\professornometext{Arthur Araruna}

\usepackage{footmisc}

\usepackage{fancyhdr}

  \pagestyle{fancy}
  \fancyhead{}
  \lhead{\disciplinacodigotext}
  \chead{\provatext}
  \rhead{\semestretext}
  \fancyfoot{}
  \fancyfoot[LE,RO]{\colorbox{black!20}{\thepage}}

% * NOTE: Redefinição feita para evitar confusão com exponenciações.
\renewcommand{\thefootnote}{(\roman{footnote})}

\begin{document}
\noindent
\begin{minipage}{\textwidth}
	\begin{tabularx}{\textwidth}{cXc}
		\multirow{4}{*}{\includegraphics[height=1.8cm]{../brasaoufc.pdf}} & {\bf UNIVERSIDADE FEDERAL DO CEARÁ}         & \multirow{3}{*}{} \\
		                                                                  & {\bf Campus de Quixadá}                     &                   \\
		                                                                  & \professorabbrv\enspace\professornometext   &                   \\
		                                                                  & \disciplinacodigotext - \disciplinanometext & {}                \\
	\end{tabularx}
\end{minipage}

\part*{\provatitle}

\section{Conceitos}

Vamos relembrar o nossos conceitos que serão necessários para o resultado a ser demonstrado:

\begin{definition}
	\emph{Rotação simples} (à direita ou à esquerda) é uma das transformações estruturais representadas a seguir, realizadas sobre a raiz de uma sub-árvore.

	\begin{figure}[h]
		\centering
		\caption{\label{fig:rot-simples}}
		\begin{tikzpicture}[transform shape, scale=0.8]
	\begin{scope}[nodes={draw, circle, minimum width=2.5ex}, sibling distance=5em, level distance=5ex]
		% Sub-árvore que representa *antes* da rotação simples à direita
		\draw node (p pre) {\( p \)}
		child {
				node (u pre) {\( u \)}
				child [child anchor=apex] {
						node [draw, isosceles triangle, shape border rotate=90, anchor=apex, minimum width=3em] (T1 pre) {\( T_1 \)}
					}
				child [child anchor=apex] {
						node [draw, isosceles triangle, shape border rotate=90, anchor=apex, minimum width=3em] (T2 pre) {\( T_2 \)}
					}
			}
		child [child anchor=apex] {
				node [draw, isosceles triangle, shape border rotate=90, anchor=apex, minimum width=3em] (T3 pre) {\( T_3 \) }
			}
		;
	\end{scope}

	% Arco que chama atenção para operação de rotação ser à direita na raiz
	\draw [->, red!80] ([shift=(180:1.5em)]p pre.center) arc (180:20:1.5em);

	\begin{scope}[nodes={draw, circle, minimum width=2.5ex}, sibling distance=5em, level distance=5ex]
		% Sub-árvore que representa *antes* de rotação simples à esquerda
		\draw node [right=16em of p pre] (u post) {\( u \)}
		child [child anchor=apex] {
				node [draw, isosceles triangle, shape border rotate=90, anchor=apex, minimum width=3em] (T1 post) {\( T_1 \) }
			}
		child {
				node (p post) {\( p \)}
				child [child anchor=apex] {
						node [draw, isosceles triangle, shape border rotate=90, anchor=apex, minimum width=3em] (T2 post) {\( T_2 \)}
					}
				child [child anchor=apex] {
						node [draw, isosceles triangle, shape border rotate=90, anchor=apex, minimum width=3em] (T3 post) {\( T_3 \)}
					}
			}
		;
	\end{scope}

	% Arco que chama atenção para operação de rotação ser à direita na raiz
	\draw [->, blue] ([shift=(0:1.5em)]u post.center) arc (0:160:1.5em);

	% Setas e legenda que destacam estado *após* cada rotação
	\begin{scope}
		\coordinate (center) at ($(current bounding box.north)!0.5!(current bounding box.south)$);
		\coordinate (rr begin) at ($(center)+(-3em,2em)$);
		\coordinate (lr begin) at ($(center)+(3em,-2em)$);

		\draw [-{Latex[fill=none]}, red!80, double, double distance=1.5pt] (rr begin) --  +(0:6em) node [draw=none, midway, label={above:{à direita (em \( p \))}}] {};
		\draw [-{Latex[fill=none]}, blue, double, double distance=1.5pt] (lr begin) -- +(180:6em) node [draw=none, midway, label={below:{à esquerda (em \( u \))}}] {};
	\end{scope}
\end{tikzpicture}
	\end{figure}
\end{definition}

\begin{definition}
	\emph{Rotação dupla} (à direita ou à esquerda) é uma das transformações estruturais representadas a seguir, realizadas sobre a raiz de uma sub-árvore.

	\begin{figure}[h]
		\centering
		\caption{\label{fig:rot-duplas}}
		\begin{tikzpicture}[transform shape, scale=0.8]
	\begin{scope}[nodes={draw, circle, minimum width=2.5ex}, sibling distance=5em, level distance=5ex]
		% Sub-árvore que representa *antes* da rotação dupla à direita
		\draw node (p drr) {\( p \)}
		child [xshift=-0.5em] {
				node (u drr) {\( u \)}
				child [child anchor=apex] {
						node [draw, isosceles triangle, shape border rotate=90, anchor=apex, minimum width=3em] (T1 drr) {\( T_1 \)}
					}
				child {
						node (v drr) {\( v \)}
						child [child anchor=apex]{
								node [draw, isosceles triangle, shape border rotate=90, anchor=apex, minimum width=3em] (T2 drr) {\( T_2 \)}}
						child [child anchor=apex]{
								node [draw, isosceles triangle, shape border rotate=90, anchor=apex, minimum width=3em] (T3 drr) {\( T_3 \)}}
					}
			}
		child [xshift=0.5em, child anchor=apex] {
				node [draw, isosceles triangle, shape border rotate=90, anchor=apex, minimum width=3em] (T4 drr) {\( T_4 \) }
			}
		;
	\end{scope}

	% Arco que chama atenção para operação de rotação ser à direita na raiz
	\draw [-{>>[flex']}, red!80] ([shift=(180:1.5em)]p drr.center) arc (180:20:1.5em);


	\begin{scope}[nodes={draw, circle, minimum width=2.5ex}, sibling distance=3.5em, level distance=5ex]
		\draw node [right=17em of p drr] (v post) {\( v \)}
		child [xshift=-1.8em] {
				node (u post) {\( u \)}
				child [child anchor=apex] {
						node [draw, isosceles triangle, shape border rotate=90, anchor=apex, minimum width=3em] (T1 post) {\( T_1 \) }
					}
				child [child anchor=apex] {
						node [draw, isosceles triangle, shape border rotate=90, anchor=apex, minimum width=3em] (T2 post) {\( T_2 \) }
					}
			}
		child [xshift=1.8em] {
				node (p post) {\( p \)}
				child [child anchor=apex] {
						node [draw, isosceles triangle, shape border rotate=90, anchor=apex, minimum width=3em] (T3 post) {\( T_3 \) }
					}
				child [child anchor=apex] {
						node [draw, isosceles triangle, shape border rotate=90, anchor=apex, minimum width=3em] (T4 post) {\( T_4 \) }
					}
			}
		;
	\end{scope}


	\begin{scope}[nodes={draw, circle, minimum width=2.5ex}, sibling distance=5em, level distance=5ex]
		% Sub-árvore que representa *antes* de rotação dupla à esquerda
		\draw node [right=17em of v post] (u dlr) {\( u \)}
		child [xshift=-0.5em, child anchor=apex] {
				node [draw, isosceles triangle, shape border rotate=90, anchor=apex, minimum width=3em] (T1 dlr) {\( T_1 \) }
			}
		child [xshift=0.5em] {
				node (p dlr) {\( p \)}
				child {
						node (v dlr) {\( v \)}
						child [child anchor=apex] {
								node [draw, isosceles triangle, shape border rotate=90, anchor=apex, minimum width=3em] (T2 dlr) {\( T_2 \)}
							}
						child [child anchor=apex] {
								node [draw, isosceles triangle, shape border rotate=90, anchor=apex, minimum width=3em] (T3 dlr) {\( T_3 \)}
							}
					}
				child [child anchor=apex] {
						node [draw, isosceles triangle, shape border rotate=90, anchor=apex, minimum width=3em] (T4 dlr) {\( T_4 \)}
					}
			}
		;
	\end{scope}

	% Arco que chama atenção para operação de rotação ser à direita na raiz
	\draw [-{>>[flex']}, blue] ([shift=(0:1.5em)]u dlr.center) arc (0:160:1.5em);

	% Setas e legenda que destacam estado *após* cada rotação
	\begin{scope}
		\coordinate (left part) at ($(T4 drr)!0.5!(u post)$);
		\coordinate (right part) at ($(p post)!0.5!(T1 dlr)$);
		\coordinate (drr begin) at ($(left part)+(180:3em)$);
		\coordinate (dlr begin) at ($(right part)+(0:3em)$);

		\draw [-{Latex[fill=none]}, red!80, double, double distance=1.5pt] (drr begin) --  +(0:5em) node [draw=none, midway, label={above:{à direita (em \( p \))}}] {};
		\draw [-{Latex[fill=none]}, blue, double, double distance=1.5pt] (dlr begin) -- +(180:5em) node [draw=none, midway, label={above:{à esquerda (em \( u \))}}] {};
	\end{scope}
\end{tikzpicture}
	\end{figure}
\end{definition}

\section{Rotações preservam propriedade ABB}

Observando as operações de rotação, podemos concluir que, se a árvore dada como entrada for binária de busca, qualquer das rotações realizada em qualquer dos seus nós preservará essa propriedade. Ou seja, após a rotação a nova estrutura obtida também será binária de busca.

\subsection{Rotações simples}

Suponha que a árvore à esquerda na~\Vref{fig:rot-simples} seja binária de busca. Isso significa que todas as relações a seguir são verdadeiras(considere que quando falamos sobre uma sub-árvore estejamos falando sobre todos os seus nós):
\begin{itemize}
	\begin{multicols}{2}
		\item \( T_1 < u \)
		\item \( u < T_2 \)
		\item \( u < p \)
		\item \( p < T_3 \)
	\end{multicols}
\end{itemize}

E, para que a estrutura após uma rotação à direita (árvore da direita) mantenha a propriedade, precisamos concluir que:
\begin{itemize}
	\begin{multicols}{2}
		\item \( T_1 < u \)
		\item \( u < p \)
		\item \( T_2 < p \)
		\item \( p < T_3 \)
	\end{multicols}
\end{itemize}

Mas observe que todas as novas relações podem ser obtidas das anteriores, lembrando que a relação de ``\( < \)'' é transitiva.

Como os nós na sub-árvore rotacionada mudaram apenas de posição, mas permanecem os mesmos, essa alteração não tem como ter prejudicado a propriedade globalmente na árvore que a contém.

Por fim, perceba que o caso da rotação simples à esquerda é simétrico ao que acabamos de observar.

\subsection{Rotações duplas}

Considere agora que a árvore mais à esquerda na~\Vref{fig:rot-duplas} seja binária de busca. Isso significa que todas as relações a seguir são verdadeiras(considere que quando falamos sobre uma sub-árvore estejamos falando sobre todos os seus nós):
\begin{itemize}
	\begin{multicols}{3}
		\item \( T_1 < u \)
		\item \( u < v \)
		\item \( T_2 < v \)
		\item \( v < T_3 \)
		\item \( u < p \)
		\item \( p < T_4 \)
	\end{multicols}
\end{itemize}

E, para que a estrutura após uma rotação dupla à direita (árvore ao centro) mantenha a propriedade, precisamos concluir que:
\begin{itemize}
	\begin{multicols}{3}
		\item \( T_1 < u \)
		\item \( u < T_2 \)
		\item \( u < v \)
		\item \( v < p \)
		\item \( T_3 < p \)
		\item \( p < T_4 \)
	\end{multicols}
\end{itemize}

Da mesma forma, observe que todas as novas relações podem ser obtidas das anteriores e globalmente a situação não é alterada. O caso à esquerda também é simétrico (considere a árvore mais à direita).

\section{Rotações corrigem desbalanceamento}

Nosso problema atual consiste em corrigir situações de desbalanceamento que possam ser criadas após uma inserção em uma árvore AVL. Vamos considerar os casos gerados nessa situação e observar que podemos resolvê-los usando alguma das rotações que definimos.

Considere uma árvore \( T \) inicialmente AVL no momento imediatamente após a inserção de um novo nó \( n \). Seja \( p \) o nó mais profundo\footnote{O de maior nível dentre todos os nós.} de \( T \) que se tornou desbalanceado após essa inserção.

Primeiramente, observe que \( p \) não pode ser uma folha\footnote{O balanço de qualquer folha é sempre 0, pois ambas as suas sub-árvores são vazias.} e que ele precisa ser um ascendente de \( n \)\footnote{Caso contrário, como a inserção o teria desbalanceado, se ela não teria ocorrido em uma de suas sub-árvores?}. Por fim, a raiz da sub-árvore de \( p \) afetada pela inserção de \( n \) não pode ser o próprio \( n \)\footnote{Pois antes ela era uma sub-árvore vazia, o que não pode ter tornado \( p \) desbalanceado.}. Chamaremos, então, essa raiz de \( u \).

Existem alguns casos possíveis para os balanços e a estrutura de \( T \) nessa situação.

\subsection{Inserção na sub-árvore esquerda de \( p \)}

\subsubsection{Inserção na sub-árvore esquerda de \( u \)}

Nesse caso, certamente a configuração da sub-árvore de \( p \) e seus balanços é a representada na ~\Vref{fig:ins-rot-dir}.
\begin{figure}[hbt]
	\caption{\label{fig:ins-rot-dir} Nó \( n \) faz parte da sub-árvore \( T_1 \).}
	\centering
	\begin{tikzpicture}[transform shape, scale=0.8]
	\begin{scope}[nodes={draw, circle, minimum width=2.5ex}, sibling distance=5em, level distance=5ex]
		% Sub-árvore que representa *após* inserção gerando caso de rotação simples à direita
		\draw node (p pre) [label={[rectangle, fill=black!30]south:{\( -2 \)}}] {\( p \)}
		child {
		node (u pre) [label={[rectangle, fill=black!30]south:{\( -1 \)}}] {\( u \)}
		child [child anchor=apex] {
				node [draw, isosceles triangle, shape border rotate=90, anchor=apex, minimum width=3em] (T1 pre) {\( T_1 \)}
			}
		child [child anchor=apex] {
				node [draw, isosceles triangle, shape border rotate=90, anchor=apex, minimum width=3em] (T2 pre) {\( T_2 \)}
			}
		}
		child [child anchor=apex] {
				node [draw, isosceles triangle, shape border rotate=90, anchor=apex, minimum width=3em] (T3 pre) {\( T_3 \) }
			}
		;
	\end{scope}
\end{tikzpicture}
\end{figure}

Assim, a seguintes relações podem se estabelecidas:
\begin{enumerate}[label={\Roman*.}, series=alturas-simples]
	\item \label{item:h-T3} \( bal(p) = -2 \implies h(T_3) - h(u) = -2 \implies h(T_3) = h(u) - 2 \)
	\item \label{item:h-T2} \( bal(u) = -1 \implies h(T_2) - h(T_1) = -1 \implies h(T_2) = h(T_1) - 1 \)
	\item \label{item:h-u} \( h(u) = h(T_1) + 1 \)\footnote{Pela definição recursiva de altura e porque o balanço nos informa a sub-árvore mais alta.}
\end{enumerate}

Se substituirmos~\ref{item:h-u} em~\ref{item:h-T3}, e usarmos a expressão obtida em~\ref{item:h-T2}, encontraremos também que
\begin{enumerate}[label={\Roman*.}, resume=alturas-simples]
	\item \( h(T_3) = h(T_1) + 1 - 2 \implies h(T_3) = h(T_1) - 1 \)
	\item \label{item:h-T2-h-T3} \( h(T_2) = (h(T_3) + 1) - 1 \implies h(T_2) = h(T_3) \)
\end{enumerate}

Isso nos permite afirmar que após realizarmos uma rotação à direita em \( p \) (\Vref{fig:ins-rot-dir-post}) os novos balanços de \( p \) e \( u \) serão:
\begin{itemize}
	\item \( bal(p) = b_p = h(T_3) - h(T_2) = 0 \)\footnote{Observando~\ref{item:h-T2-h-T3}}
	\item \( bal(u) = b_u = h(p) - h(T_1) = (h(T_2) + 1) - h(T_1) = (h(T_1) - 1) + 1 - h(T_1) = 0 \)\footnote{Observando a definição recursiva de altura para \( p \) e o item~\ref{item:h-T2}}
\end{itemize}

\begin{figure}[hbt]
	\caption{\label{fig:ins-rot-dir-post}}
	\centering
	\begin{tikzpicture}[transform shape, scale=0.8]
	\begin{scope}[nodes={draw, circle, minimum width=2.5ex}, sibling distance=5em, level distance=5ex]
		% Sub-árvore que representa *após* inserção gerando caso de rotação simples à direita
		\draw node (u post) [label={[rectangle, fill=red!30]south:{\( b_u = 0 \)}}] {\( u \)}
		child [child anchor=apex] {
				node [draw, isosceles triangle, shape border rotate=90, anchor=apex, minimum width=3em] (T1 post) {\( T_1 \) }
			}
		child {
		node (p post) [label={[rectangle, fill=red!30]south:{\( b_p = 0 \)}}] {\( p \)}
		child [child anchor=apex] {
				node [draw, isosceles triangle, shape border rotate=90, anchor=apex, minimum width=3em] (T2 post) {\( T_2 \)}
			}
		child [child anchor=apex] {
				node [draw, isosceles triangle, shape border rotate=90, anchor=apex, minimum width=3em] (T3 post) {\( T_3 \)}
			}
		}
		;
	\end{scope}
\end{tikzpicture}
\end{figure}

Dessas informações, podemos concluir que a sub-árvore rotacionada passou a estar toda balanceada, já que os novos balanços são 0 e não alteramos as estruturas das sub-árvores \( T_1, T_2 \) e \( T_3 \). Além disso, podemos concluir que a altura dessa sub-árvore sob a ótica dos seus ancestrais não mudou entre antes e depois da inserção\footnote{Veja~\Vref{block:altura-nao-muda}.}.

\section{Conclusão}


\part*{Apêndice}

\section{Altura não muda após rotação}
\label{block:altura-nao-muda}


%\printbibliography

\end{document}